\documentclass[
  headings=optiontotocandhead,
	oneside,
	numbers=noenddot,
	toc=flat,
	11pt,
	parskip=full,
	listof=flat,
  numbers=noenddot,
	bibliography=totoc,
	a4paper,
  DIV=14,
  headheight=5mm,
  headinclude=true,
]{scrartcl}

% \usepackage{fontspec}
% \usepackage{titlesec}
\usepackage{float}
\usepackage{array}
\usepackage[main=ngerman]{babel} %texlive-langgerman muss installiert sein
\usepackage{graphicx}
\usepackage{subcaption}
\usepackage{xcolor}
\usepackage{relsize}
\usepackage{lastpage}

\input{glyphtounicode}
\pdfgentounicode=1

\clubpenalty = 1000 % Schusterjungen verhindern
\widowpenalty = 1000 % Hurenkinder verhindern
\displaywidowpenalty = 1000

\newcommand{\kopfbild}{\voffset7mm\includegraphics[width=25mm]{TU_Logo_RGB}}
\newcommand{\kopftitel}{\sffamily{\textbf{\large{Laborprotokoll Interferometrie}}}}

\usepackage[automark,footsepline,plainfootsepline]{scrlayer-scrpage}
\setkomafont{pageheadfoot}{\normalcolor\footnotesize\scshape}
\setkomafont{pagenumber}{\normalfont\normalsize}
\clearpairofpagestyles
\ihead{\headmark}
\ihead{\kopfbild}
\ohead{\kopftitel}
\ifoot{\smaller{David Koch, Paul Vetter}}
\ofoot{Seite \pagemark/\pageref{LastPage}}

\title{Laborprotokoll Interferometrie}
\author{David Koch, Paul Vetter}
\date{}


\begin{document}
\maketitle
\tableofcontents

\section{Theorie}
Laser? Interferenz? Koheränz Länge?
\section{Aufbau}

\section{Anhang/Messergebnisse}

Bei der interferometrischen Längenmessung wurden 5 Messungen durchgeführt und anschließend wurden die Ergebnisse, also die Anzahl der Interferenzringe für eine bestimmte zurückgelegte Distanz des SP2, gemittelt.

\begin{table}[H]
\centering
\setlength{\tabcolsep}{10pt}
\begin{tabular}{ |c|c| }
  \multicolumn{1}{c}{Distanz [$\mu m$]}  & \multicolumn{1}{c}{$\sum$ Ringe} \\
  40 & 127 \\
  30 & 87 \\
  10 & 31 \\
  50 & 162 \\
  40 & 133 \\
\end{tabular}
\caption{Interferometrische Wellenlängenmessung}
\label{tab:1}
\end{table}

Die Messungen wurden aufgrund von zeitlichen Einschränkungen jeweils für unterschiedliche (kürzere) Distanzen und nicht immer eine vollständige Umdrehung der Feinverstellung der Differential-Schraube durchgeführt. Dies ändert jedoch nichts am Ergebnis, denn mit der gemittelten Anzahl der Interferenzringe lässt sich auf die Wellenlänge des Laserlichts zurückführen:

$  $ 

INSERT RECHNUNG HERE

\begin{table}[H]
\centering
\setlength{\tabcolsep}{15pt}
\begin{tabular}{ |c|c| }
  \multicolumn{1}{c}{Distanz [$\mu m$]}  & \multicolumn{1}{c}{$\sum$ Ringe} \\
  13 & 50 \\
  12 & 50 \\
\end{tabular}
\caption{Natrium Wellenlängemessung}
\label{tab:2}
\end{table}

8. $d = 8 mm$

\end{document}

%%% Local Variables:
%%% mode: LaTeX
%%% TeX-master: t
%%% End:
